\providecommand{\verfasser}{Jannes Bantje}

%--Konfiguration von scrheadings
\setheadsepline{1pt}[\color{light_gray}]
\pagestyle{scrheadings}
%\fancyhf[H, F]{}
\clearscrheadfoot

\providecommand{\shortFach}{\fach}
\lehead{ \includegraphics[height=0.6 cm,keepaspectratio]{../!config/Bilder/Logo_WWU_Muenster_light_gray.pdf}}
\rehead{\rule{0cm}{0.6cm}\footnotesize \sffamily \color{light_gray} \verfasser -- Mitschrift \shortFach}
\lohead{\rule{0cm}{0.6cm} \footnotesize \sffamily \color{light_gray} Stand: \today \; \thistime[:]}
\rohead{\includegraphics[height=0.6 cm,keepaspectratio]{../!config/Bilder/fb10logo_gray.pdf}}


\ofoot[{ \color{dark_gray} \LARGE \sffamily \thepage}]{{ \color{dark_gray} \LARGE \sffamily \thepage}} %hier wir auch der plain Stil bearbeitet!
\automark{section}
\ifoot{ \color{dark_gray} \small \leftmark}

%--Metadaten
\providecommand{\mail}{j.bantje@wwu.de}
\author{\verfasser}
\titlehead{\includegraphics[height=1.5cm, keepaspectratio]{../!config/Bilder/Logo_WWU_Muenster.pdf}%
\hfill \includegraphics[height=1.3cm, keepaspectratio]{../!config/Bilder/fb10logo.pdf}}
\title{Skript \fach}
\subtitle{Mitschrift der Vorlesung  \enquote{\fach} von \prof}

\publishers{
Aktuelle Version verfügbar bei: \bigskip\\
\normalsize
\newcommand{\dieBreite}{11cm}
\begin{minipage}{4cm}
	\qrcode[height=3.3cm, version=6]{https://github.com/JaMeZ-B/latex-wwu}
\end{minipage}
\hfill
\begin{minipage}{\dieBreite}
	\includegraphics[height=0.6cm, keepaspectratio]{../!config/Bilder/github_octo.pdf}
	\includegraphics[height=0.6cm, keepaspectratio]{../!config/Bilder/GitHub_Logo.pdf}\\
	\url{https://github.com/JaMeZ-B/latex-wwu} \smallskip\\
	GitHub ist eine Internetplattform, auf der viele OpenSource-Projekte gehostet werden. Diese Plattform nutzen wir zur Zusammenarbeit, also findet man hier neben den PDFs auch die 
	\TeX-Dateien. Außerdem ist über diese Plattform auch direktes Mitarbeiten möglich, siehe nächste Seite.
\end{minipage}\\[1cm]
\begin{minipage}{4cm}
	\qrcode[height=3.3cm, version=6]{https://uni-muenster.sciebo.de/public.php?service=files&t=965ae79080a473eb5b6d927d7d8b0462}
\end{minipage}
\hfill
\begin{minipage}{\dieBreite}
	\raisebox{-2pt}{\includegraphics[height=0.6cm, keepaspectratio]{../!config/Bilder/sciebo_logo.pdf}}
	\resizebox{!}{0.5cm}{\large \textbf{sciebo}} {\large die Campuscloud} \\
	\resizebox{\dieBreite}{!}{\footnotesize\url{https://uni-muenster.sciebo.de/public.php?service=files&t=965ae79080a473eb5b6d927d7d8b0462}}\smallskip\\
	Sciebo ist ein Dropbox-Ersatz der Hochschulen in NRW, der von der Uni Münster in leitender Position auf Basis der OpenSource-Software Owncloud aufgebaut wurde. Wenn man auf den 
	Link klickt, kann man die Freigabe zum eigenen Speicher hinzufügen und hat dann immer automatisch die aktuellste Version.
\end{minipage}\\[1cm]
\begin{minipage}{4cm}
	\qrcode[height=3.3cm, version=6]{btsync://B6WH2DISQ5QVYIRYIEZSF4ZR2IDVKPN3I?n=latex_share}
\end{minipage}
\hfill
\begin{minipage}{\dieBreite}
	\raisebox{-2pt}{\includegraphics[height=0.6cm, keepaspectratio]{../!config/Bilder/bt_sync_logo.pdf}}
	\resizebox{!}{0.5cm}{\large \textbf{Bittorrent} Sync}\\
	\texttt{B6WH2DISQ5QVYIRYIEZSF4ZR2IDVKPN3I} \smallskip\\
	BTSync ist ein peer-to-peer Dateisynchronisations-Tool. Dabei werden die Dateien nur auf den Computern der Teilnehmer an einer Freigabe gespeichert. Ein Mini-Computer ist permanent 
	online, sodass jederzeit die aktuellste Version verfügbar ist. \hrefsym{http://www.getsync.com/intl/de/}{Clients} gibt es für jedes Betriebssystem.
	Zugang ist über das obige \enquote{Secret} bzw. den QR-Code möglich
\end{minipage}
% \begin{minipage}[t]{0.45\textwidth}
% 	\centering{
% 	\qrcode[height=3.5cm, version=4]{https://github.com/JaMeZ-B/latex-wwu}
% 	\medskip\\
% 	\includegraphics[height=0.4cm, keepaspectratio]{../!config/Bilder/github_octo.pdf}
% 	\includegraphics[height=0.4cm, keepaspectratio]{../!config/Bilder/GitHub_Logo.pdf} (inklusive Sourcecode)\\
% 	{\small\url{https://github.com/JaMeZ-B/latex-wwu}}}
% \end{minipage}
% \hfill
% \begin{minipage}[t]{0.45\textwidth}
% 	\centering{
% 	\qrcode[height=3.5cm, version=4]{btsync://B6WH2DISQ5QVYIRYIEZSF4ZR2IDVKPN3I?n=latex_share}
% 	\medskip\\
% 	\includegraphics[height=0.4cm, keepaspectratio]{../!config/Bilder/bt_sync_logo.pdf}
% 	{\large \textbf{Bittorrent} Sync} \\
% 	\texttt{B6WH2DISQ5QVYIRYIEZSF4ZR2IDVKPN3I}
% 	}
% \end{minipage} \bigskip\\
% \begin{minipage}[t]{0.45\textwidth}
% 	\centering{
% 	\qrcode[height=3.5cm, version=4]{https://uni-muenster.sciebo.de/public.php?service=files&t=965ae79080a473eb5b6d927d7d8b0462}
% 	\medskip\\
% 	\includegraphics[height=0.4cm, keepaspectratio]{../!config/Bilder/sciebo_logo.pdf}
% 	{\large \textbf{sciebo} die Campuscloud} \\
% 	\hrefsym{https://uni-muenster.sciebo.de/public.php?service=files&t=965ae79080a473eb5b6d927d7d8b0462}{\texttt{LaTeX-Freigabe}}
% 	}
% \end{minipage}
}
